% Generated draft: requires mathematical and source review.
\documentclass[11pt]{article}
\usepackage[margin=25mm]{geometry}
\usepackage{fontspec}
\setmainfont{Latin Modern Roman}
\usepackage{amsmath,amssymb,mathtools,unicode-math}
\setmathfont{Latin Modern Math}
\usepackage{enumitem,needspace,xurl,hyperref}
\setlength{\parindent}{0pt}
\setlength{\parskip}{5pt}
\setlength{\emergencystretch}{3em}
\newcommand{\N}{\mathbb{N}}
\newcommand{\Z}{\mathbb{Z}}
\newcommand{\Q}{\mathbb{Q}}
\newcommand{\R}{\mathbb{R}}
\newcommand{\C}{\mathbb{C}}
\newcommand{\F}{\mathbb{F}}
\newcommand{\A}{\mathbb{A}}
\newcommand{\PP}{\mathbb{P}}
\newcommand{\E}{\mathbb{E}}
\newcommand{\eps}{\varepsilon}
\newcommand{\dd}{\,\mathrm{d}}
\begin{document}
\section*{UnsolvedMath ID: 701125}
\textbf{Source problem code:} catalogue entry
\subsection{Definitions and mathematical statement}
For a subset \(A\subseteq G\), write
\[
C_G(A)=\{g\in G:ga=ag\text{ for every }a\in A\}.
\]
Say that \(G\) satisfies the minimum condition for centralizers if every descending sequence of centralizers
\[
C_G(A_1)\supseteq C_G(A_2)\supseteq\cdots
\]
eventually stabilizes. A subset \(X\subseteq G\) is normal if \(g^{-1}Xg=X\) for every \(g\in G\). Use the left-normed commutator convention
\[
[a,b]=a^{-1}b^{-1}ab,\qquad [a,{}_1b]=[a,b],\qquad [a,{}_{n+1}b]=[[a,{}_nb],b].
\]
A group is locally nilpotent if every finitely generated subgroup is nilpotent. A group \(H\) is hypercentral if \(Z_\alpha(H)=H\) for some ordinal \(\alpha\), where \(Z_0(H)=1\), \(Z_{\alpha+1}(H)\) is the inverse image of \(Z(H/Z_\alpha(H))\), and \(Z_\lambda(H)=\bigcup_{\alpha<\lambda}Z_\alpha(H)\) for every limit ordinal \(\lambda\).

\textbf{Problem (F. O. Wagner, 2002).} Let \(G\) satisfy the minimum condition for centralizers, and let \(X\subseteq G\) be a normal subset. Suppose there exists a function
\[
f:X\times X\longrightarrow\Z_{\geq 1}
\]
such that
\[
[x,{}_{f(x,y)}y]=1\qquad\text{for every }x,y\in X.
\]
Must \(H=\langle X\rangle\) be locally nilpotent? The source further describes the desired conclusion as ``locally nilpotent (whence hypercentral)''; this parenthetical is retained under the stated ambient hypotheses and is not an assertion about arbitrary locally nilpotent groups.

No uniform bound on \(f(x,y)\) is assumed. The source reports an affirmative answer when there exists \(N\in\Z_{\geq1}\) such that \(f(x,y)\leq N\) for every \(x,y\in X\), attributing this bounded case to F. O. Wagner, \emph{J. Algebra} \textbf{217}, no. 2 (1999), 448--460.
\subsection{Short English statement}
In a group with the minimum condition for centralizers, suppose a conjugation-invariant subset has the property that repeatedly commuting any one of its elements with another eventually gives the identity, with the number of repetitions allowed to depend on the ordered pair. Must the subset generate a locally nilpotent subgroup, and hence a hypercentral subgroup as stated in the source? The source reports that a uniform bound on the repetitions gives an affirmative answer.
\subsection{Sources}
\begin{itemize}
\item Local UnsolvedMath registry, permanent ID \(701125\), exact source code \texttt{catalogue entry}, title \emph{Kourovka Notebook Problem 15.37}; supplied upstream revision \texttt{d371ee917ac9647fbb06cd6fc529deec6f3dda82}. The supplied status is \texttt{open}; this is a registry label, not an independently verified assessment.
\item \emph{The Kourovka Notebook (Unsolved Problems in Group Theory)}, problem 15.37, proposed by F. O. Wagner in the 15th issue (2002); supplied reference to the 21st issue edition (2026), arXiv:1401.0300v48. The supplied collection description names E. I. Khukhro and V. D. Mazurov as editors.
\item F. O. Wagner, \emph{J. Algebra} \textbf{217}, no. 2 (1999), 448--460. The supplied statement cites this reference for the affirmative bounded case; no article title is supplied.
\end{itemize}
\end{document}
